\documentclass[11pt]{article}
\usepackage[a4paper,margin=2.5cm]{geometry}
\usepackage{amsmath}
\usepackage{amssymb}
\usepackage{enumitem}
\usepackage{parskip}
\setlist[enumerate,2]{label=(\alph*), itemsep=1pt, topsep=2pt}
\setlength{\parindent}{0pt}
\newcommand{\correct}[1]{\textbf{#1}\;$\checkmark$}
\begin{document}
\section*{Moran process}
\begin{enumerate}[itemsep=6pt]
  \item At each step of a Moran process ...
  \begin{enumerate}
    \item every individual reproduces once
    \item the population size grows by one
    \item the least fit individual always dies
    \item one individual reproduces (proportional to fitness) and one is chosen uniformly to die
  \end{enumerate}
  \item The fixation probability of a type is ...
  \begin{enumerate}
    \item the probability the population eventually consists entirely of that type
    \item the probability the type dies out
    \item the fraction of the population it starts with
    \item its fitness divided by the average fitness
  \end{enumerate}
  \item Under neutral drift (a mutant of the same fitness as residents), a single mutant fixes with probability ...
  \begin{enumerate}
    \item \(1/N\)
    \item \(0\)
    \item \(1/2\)
    \item \(1\)
  \end{enumerate}
  \item A Moran process with two types has ...
  \begin{enumerate}
    \item two absorbing states: all of one type or all of the other
    \item no absorbing states
    \item an absorbing state at every population split
    \item exactly one absorbing state
  \end{enumerate}
  \item For an advantageous mutant with relative fitness \(r > 1\), as \(N \to \infty\) the fixation probability tends to ...
  \begin{enumerate}
    \item \(1 - 1/r\), which is bounded away from 1
    \item \(1\)
    \item \(1/N\)
    \item \(0\)
  \end{enumerate}
  \item The relative fitness \(r\) of a mutant is ...
  \begin{enumerate}
    \item the size of the population
    \item the number of mutants
    \item the probability it dies
    \item its fitness divided by the resident fitness
  \end{enumerate}
  \item In a game-based Moran process, the fitness of an individual depends on ...
  \begin{enumerate}
    \item only its own type
    \item the composition of the whole population (its payoff in the current state)
    \item nothing; all fitnesses are equal
    \item the round number
  \end{enumerate}
  \item Selection in the Moran process acts through ...
  \begin{enumerate}
    \item the population size
    \item the birth step: fitter individuals are more likely to reproduce
    \item neither step
    \item the death step: fitter individuals are more likely to die
  \end{enumerate}
  \item The death step of the Moran process chooses an individual ...
  \begin{enumerate}
    \item that just reproduced
    \item uniformly at random
    \item with the lowest fitness
    \item proportional to fitness
  \end{enumerate}
  \item The two-type fixation formula uses the ratios \(\gamma_i\) defined as ...
  \begin{enumerate}
    \item the fitness of type 1
    \item the population size
    \item \(p_{i,i-1} / p_{i,i+1}\), the down-step over up-step probabilities
    \item \(1/N\)
  \end{enumerate}
  \item Under strong selection (\(w = 1\)) ...
  \begin{enumerate}
    \item the process never fixes
    \item fitness equals the raw payoff, so payoff differences act at full strength
    \item all fitnesses are equal
    \item the population grows
  \end{enumerate}
  \item The boundary conditions for the fixation probability are ...
  \begin{enumerate}
    \item \(\rho_0 = \rho_N = 1/N\)
    \item \(\rho_0 = 1\) and \(\rho_N = 0\)
    \item \(\rho_0 = \rho_N = 1/2\)
    \item \(\rho_0 = 0\) and \(\rho_N = 1\)
  \end{enumerate}
  \item A single mutant is favoured by selection when its fixation probability satisfies ...
  \begin{enumerate}
    \item \(\rho_1 > 1/N\)
    \item \(\rho_1 = 0\)
    \item \(\rho_1 = 1\)
    \item \(\rho_1 < 1/N\)
  \end{enumerate}
  \item The Moran process is repeated until ...
  \begin{enumerate}
    \item a fixed number of steps have passed
    \item the population doubles
    \item every individual has reproduced once
    \item only one type remains in the population
  \end{enumerate}
  \item A disadvantageous mutant (\(r < 1\)) has a fixation probability that is ...
  \begin{enumerate}
    \item exactly \(1\)
    \item above the neutral value \(1/N\)
    \item exactly \(1/2\)
    \item below the neutral value \(1/N\)
  \end{enumerate}
  \item The two-type fixation probability is \(\rho_i = \dfrac{1 + \sum_{j=1}^{i-1}\prod_{k=1}^{j}\gamma_k}{1 + \sum_{j=1}^{N-1}\prod_{k=1}^{j}\gamma_k}\), where \(\gamma_k\) equals ...
  \begin{enumerate}
    \item the mutant fitness alone
    \item \(1/N\)
    \item \(f_R(k)/f_M(k)\), the resident-to-mutant fitness ratio
    \item \(f_M(k)/f_R(k)\)
  \end{enumerate}
  \item In a constant-fitness process with mutant fitness \(r\) and resident fitness 1, the ratio \(\gamma_k\) is ...
  \begin{enumerate}
    \item \(r\) for every state \(k\)
    \item different in every state
    \item \(r^{-1}\) for every state \(k\)
    \item equal to \(1/N\)
  \end{enumerate}
  \item Weak selection in the Moran process corresponds to a selection intensity \(w\) that is ...
  \begin{enumerate}
    \item close to 0
    \item close to 1
    \item negative
    \item greater than \(N\)
  \end{enumerate}
  \item Throughout a Moran process the total population size ...
  \begin{enumerate}
    \item shrinks over time
    \item grows by one each step
    \item fluctuates randomly
    \item stays constant at \(N\), since each birth replaces a death
  \end{enumerate}
  \item Neutral drift is the special case of the Moran process in which the relative fitness is ...
  \begin{enumerate}
    \item \(r > 1\)
    \item \(r = 0\)
    \item \(r < 1\)
    \item \(r = 1\)
  \end{enumerate}
\end{enumerate}
\end{document}
